\chapter{快速上手：环境、张量与 autograd}
\section{本章定位}
PyTorch 入门面试题高频考三类能力：环境能否跑通，张量 API 是否熟练，autograd 是否理解。
本章用精炼示例覆盖安装、GPU、张量、NumPy 互转和自动求导。

\section{环境与 GPU 安装}
\subsection{pip 与 conda 安装}
PyTorch 2.x 的安装命令应以官网为准，因为系统、Python 版本、CUDA 版本都会影响安装包。
常见 pip 写法如下：
\begin{lstlisting}[style=plainstyle]
python -m pip install torch torchvision torchaudio
\end{lstlisting}
常见 conda 写法如下：
\begin{lstlisting}[style=plainstyle]
conda install pytorch torchvision torchaudio -c pytorch
\end{lstlisting}
如果要使用 NVIDIA GPU，需要显卡驱动、PyTorch CUDA 构建和运行时 GPU 可见性都正确。
很多 PyTorch wheel 已自带 CUDA 运行时库，本机不一定要额外安装完整 CUDA Toolkit。

\subsection{查版本、查 GPU、选 device}
安装后先运行：
\begin{lstlisting}
import torch
print(torch.__version__)
print(torch.version.cuda)
print(torch.cuda.is_available())
print(torch.cuda.device_count())
\end{lstlisting}
标准 device 写法如下：
\begin{lstlisting}
import torch
device = "cuda" if torch.cuda.is_available() else "cpu"
x = torch.randn(2, 3, device=device)
print(x.device)
\end{lstlisting}
\begin{trap}
\code{torch.version.cuda} 表示 PyTorch 构建对应的 CUDA 版本，不代表机器一定有可用 GPU。
运行时判断 GPU，应使用 \code{torch.cuda.is_available()}。
\end{trap}

\section{张量基础}
\subsection{创建张量}
常见创建函数包括 \code{torch.tensor}、\code{zeros}、\code{ones}、\code{randn} 和 \code{arange}：
\begin{lstlisting}
import torch
a = torch.tensor([[1, 2], [3, 4]])
z = torch.zeros(2, 3)
o = torch.ones(2, 3)
r = torch.randn(2, 3)
t = torch.arange(0, 6)
print(a)
print(z.shape, o.shape, r.dtype, t)
\end{lstlisting}
\code{torch.tensor} 会根据输入推断 dtype。训练中常用 \code{torch.float32} 表示特征和权重，分类标签常用 \code{torch.long}。

\subsection{dtype 与转换}
dtype 影响数值范围、显存占用和算子行为：
\begin{lstlisting}
import torch
x = torch.tensor([1, 2, 3])
y = x.float()
label = torch.tensor([0, 1, 2], dtype=torch.long)
half_y = y.to(torch.float16)
print(x.dtype, y.dtype, label.dtype, half_y.dtype)
\end{lstlisting}
常用转换：\code{x.float()}、\code{x.long()}、\code{x.to(dtype=torch.float16)}、\code{x.to(device)}。

\subsection{shape、索引与切片}
形状可用 \code{shape} 或 \code{size()} 查看，索引和切片接近 NumPy：
\begin{lstlisting}
import torch
x = torch.arange(12).reshape(3, 4)
print(x.shape)
print(x.size())
print(x[0])
print(x[:, 1])
print(x[1:3, 2:])
\end{lstlisting}
切片通常返回 view，可能与原张量共享底层存储。

\subsection{广播规则}
广播从最后一个维度向前比较：
\begin{enumerate}[label=\arabic*.]
  \item 维度相等，可以广播。
  \item 其中一个维度为 $1$，可以扩展。
  \item 维度更少的一方，在前面补 $1$ 再比较。
  \item 既不相等也没有 $1$，则不能广播。
\end{enumerate}
例子：
\begin{lstlisting}
import torch
x = torch.ones(2, 3)
b = torch.tensor([10.0, 20.0, 30.0])
y = x + b
print(y)
print(y.shape)
\end{lstlisting}
这里 \code{x} 是 $(2,3)$，\code{b} 是 $(3)$，可看成 $(1,3)$，再扩展到 $(2,3)$。
\code{keepdim=True} 常用于保留规约维度，方便后续广播：
\begin{lstlisting}
import torch
x = torch.randn(4, 3)
mean = x.mean(dim=0, keepdim=True)
std = x.std(dim=0, keepdim=True)
y = (x - mean) / (std + 1e-6)
print(mean.shape, y.shape)
\end{lstlisting}

\subsection{reshape、view、permute 与 transpose}
常见形状操作如下：
\begin{lstlisting}
import torch
x = torch.arange(12)
a = x.reshape(3, 4)
b = a.view(2, 6)
c = a.transpose(0, 1)
d = a.reshape(1, 3, 4).permute(0, 2, 1)
print(a.shape, b.shape, c.shape, d.shape)
\end{lstlisting}
\code{view} 要求张量 contiguous；\code{reshape} 会尽量返回 view，必要时复制。
\code{transpose} 交换两个维度，\code{permute} 可重排多个维度。
\begin{lstlisting}
import torch
x = torch.arange(12).reshape(3, 4)
y = x.t()
print(y.is_contiguous())
z = y.contiguous().view(12)
print(z)
\end{lstlisting}
\begin{trap}
\code{transpose} 和 \code{permute} 常返回非 contiguous 张量，之后直接 \code{view} 可能报错。
稳妥写法是先 \code{contiguous()} 再 \code{view}，或直接用 \code{reshape}。
\end{trap}

\subsection{CPU 到 GPU 迁移}
张量和模型必须在同一个 device 上计算：
\begin{lstlisting}
import torch
device = "cuda" if torch.cuda.is_available() else "cpu"
x = torch.randn(2, 3)
x = x.to(device)
y = torch.ones(2, 3, device=device)
z = x + y
print(z.cpu().device)
\end{lstlisting}
\code{x.to(device)} 通常返回新张量；不接返回值，原 \code{x} 可能仍在原 device 上。

\subsection{与 NumPy 互转}
CPU 张量和 NumPy 数组可以互转，且通常共享内存：
\begin{lstlisting}
import torch
x = torch.tensor([1.0, 2.0, 3.0])
a = x.numpy()
a[0] = 100.0
print(x)
y = torch.from_numpy(a)
y[1] = 200.0
print(a)
\end{lstlisting}
GPU 张量转 NumPy 前要先到 CPU，训练结果通常还要先 \code{detach()}：
\begin{lstlisting}
import torch
device = "cuda" if torch.cuda.is_available() else "cpu"
x = torch.randn(3, device=device)
a = x.detach().cpu().numpy()
print(a)
\end{lstlisting}
\begin{trap}
NumPy 互转共享内存很高效，但修改一方会影响另一方。
如果需要独立副本，可用 \code{x.detach().cpu().clone().numpy()} 或 NumPy 的 \code{copy()}。
\end{trap}

\subsection{常用算子}
矩阵乘法、规约、拼接和维度增删是高频操作：
\begin{lstlisting}
import torch
a = torch.randn(2, 3)
b = torch.randn(3, 4)
c = a @ b
d = torch.matmul(a, b)
x = torch.randn(2, 3, 4)
s = x.sum()
m = x.mean(dim=1)
v, idx = x.max(dim=2)
p = torch.randn(2, 3)
q = torch.randn(2, 3)
cat0 = torch.cat([p, q], dim=0)
stack0 = torch.stack([p, q], dim=0)
u = p.unsqueeze(0)
sq = u.squeeze(0)
print(c.shape, d.shape)
print(s.shape, m.shape, v.shape, idx.shape)
print(cat0.shape, stack0.shape, u.shape, sq.shape)
\end{lstlisting}
\code{cat} 沿已有维度拼接；\code{stack} 新增一个维度后堆叠。

\section{autograd 自动求导}
\subsection{requires grad、计算图与叶子张量}
设置 \code{requires_grad=True} 后，PyTorch 会记录可微运算并构建动态计算图：
\begin{lstlisting}
import torch
x = torch.tensor([1.0, 2.0, 3.0], requires_grad=True)
y = x * x + 2.0
loss = y.mean()
print(loss)
print(loss.grad_fn)
\end{lstlisting}
\code{x} 是叶子张量，\code{y} 和 \code{loss} 是运算产生的非叶子张量。
反向传播后，梯度通常保存在叶子张量的 \code{grad} 字段中。

\subsection{backward、grad 与梯度累加}
标量 loss 可直接调用 \code{backward()}：
\begin{lstlisting}
import torch
x = torch.tensor([1.0, 2.0, 3.0], requires_grad=True)
loss = (x * x).sum()
loss.backward()
print(x.grad)
\end{lstlisting}
这里 $loss=x_1^2+x_2^2+x_3^2$，所以 $\partial loss/\partial x_i=2x_i$。
如果输出不是标量，需要传入同形状上游梯度。
PyTorch 梯度默认累加，因此训练循环要清零：
\begin{lstlisting}
import torch
model = torch.nn.Linear(3, 1)
optimizer = torch.optim.SGD(model.parameters(), lr=0.1)
x = torch.randn(4, 3)
y = torch.randn(4, 1)
pred = model(x)
loss = torch.nn.functional.mse_loss(pred, y)
optimizer.zero_grad()
loss.backward()
optimizer.step()
\end{lstlisting}

\subsection{no grad、detach 与 in-place 风险}
\code{with torch.no_grad()} 临时关闭梯度记录，常用于验证和推理：
\begin{lstlisting}
import torch
x = torch.tensor([1.0, 2.0, 3.0], requires_grad=True)
with torch.no_grad():
    y = x * 2.0
print(y.requires_grad)
\end{lstlisting}
\code{detach()} 返回从当前计算图分离的新张量：
\begin{lstlisting}
import torch
x = torch.tensor([1.0, 2.0, 3.0], requires_grad=True)
y = x * 2.0
z = y.detach()
print(y.requires_grad)
print(z.requires_grad)
\end{lstlisting}
in-place 操作会直接修改原张量，常以末尾下划线表示，例如 \code{zero_()}。
它可能破坏 autograd 保存的中间值，因此对需要梯度的张量要谨慎。

\section{高频面试题}
\begin{interview}
\textbf{view 和 reshape 有什么区别？什么时候 view 会失败？}
\end{interview}
\begin{answer}
\code{view} 只改变形状视图，要求底层内存 contiguous；非 contiguous 张量直接 \code{view} 可能失败。
\code{reshape} 会尽量返回 view，做不到时会复制数据。一句话：\code{view} 更严格，\code{reshape} 更省心但可能有拷贝。
\end{answer}
\begin{interview}
\textbf{PyTorch 的广播规则是什么？}
\end{interview}
\begin{answer}
从最后一维向前比较。维度相等，或其中一个为 $1$，就能广播；缺失的前置维度按 $1$ 处理。
例如 $(2,3)$ 与 $(3)$ 可看成 $(2,3)$ 与 $(1,3)$，结果为 $(2,3)$。
\end{answer}
\begin{interview}
\textbf{requires\_grad、detach 和 no\_grad 有什么区别？}
\end{interview}
\begin{answer}
\code{requires_grad=True} 表示需要追踪梯度；\code{detach()} 把已有结果从当前计算图分离；
\code{torch.no_grad()} 在代码块内暂停记录新运算，常用于推理。
\end{answer}
\begin{interview}
\textbf{为什么 numpy() 与 tensor 可能共享内存？}
\end{interview}
\begin{answer}
CPU tensor 调用 \code{numpy()} 时通常不复制数据，而是与 NumPy 数组共享底层内存以提升效率。
风险是修改一方会影响另一方。需要副本时使用 \code{clone()} 或 NumPy 的 \code{copy()}。
\end{answer}
\begin{interview}
\textbf{标量 .item() 有什么作用？}
\end{interview}
\begin{answer}
\code{.item()} 把单元素 tensor 转成 Python 数值，常用于日志和打印。
它只能用于单元素张量；若张量在 GPU 上，调用 \code{.item()} 会触发 CPU 与 GPU 同步，频繁调用会拖慢训练。
\end{answer}
\begin{quickcard}\begin{itemize}
  \item 查版本：\code{torch.__version__}；查 CUDA 构建：\code{torch.version.cuda}；查 GPU：\code{torch.cuda.is_available()}；标准 device：\code{device = "cuda" if torch.cuda.is_available() else "cpu"}。
  \item 广播从后往前比；相等或有一个为 $1$ 才兼容；\code{view} 要求 contiguous，\code{reshape} 更通用但可能复制。
  \item \code{cat} 沿已有维度拼接；\code{stack} 新增维度堆叠；CPU tensor 与 NumPy 通常共享内存。
  \item \code{requires_grad=True} 追踪梯度；梯度默认累加，训练循环中要 \code{optimizer.zero_grad()}。
  \item \code{torch.no_grad()} 关闭记录；\code{detach()} 分离已有计算图；\code{.item()} 会取 Python 标量并可能触发 GPU 同步。
\end{itemize}\end{quickcard}
